\documentclass[11pt,a4paper]{article}
% ---- Préambule commun ----
\usepackage[utf8]{inputenc}
\usepackage[T1]{fontenc}
\usepackage{lmodern}
\usepackage[french]{babel}
\usepackage[a4paper,margin=2.2cm,top=2.6cm,bottom=2.4cm]{geometry}
\usepackage{amsmath,amssymb,amsthm}
\usepackage{enumitem}
\usepackage{xcolor}
\usepackage{listings}
\usepackage{fancyhdr}
\usepackage[most]{tcolorbox}
\usepackage{hyperref}

\definecolor{teal}{HTML}{026573}
\definecolor{tealclair}{HTML}{E6F2F3}
\definecolor{grisclair}{HTML}{F5F5F5}
\definecolor{orangefonce}{HTML}{B35900}

\hypersetup{colorlinks=true,linkcolor=teal,urlcolor=teal,citecolor=teal}

\setlength{\parindent}{0pt}
\setlength{\parskip}{3pt}

% ---- Style Python ----
\lstdefinestyle{python}{
  language=Python,
  basicstyle=\ttfamily\small,
  keywordstyle=\color{teal}\bfseries,
  commentstyle=\color{gray}\itshape,
  stringstyle=\color{orangefonce},
  showstringspaces=false,
  numbers=left,
  numberstyle=\tiny\color{gray},
  numbersep=8pt,
  frame=leftline,
  framerule=1.2pt,
  rulecolor=\color{teal},
  xleftmargin=14pt,
  backgroundcolor=\color{grisclair},
  breaklines=true,
  literate={é}{{\'e}}1 {è}{{\`e}}1 {ê}{{\^e}}1 {à}{{\`a}}1 {ô}{{\^o}}1 {û}{{\^u}}1 {î}{{\^i}}1 {ç}{{\c c}}1
}
\lstset{style=python}

% ---- Compteur d'exercices ----
\newcounter{exo}
\newcommand{\exo}[1]{%
  \stepcounter{exo}%
  \vspace{10pt}%
  \noindent\begin{tcolorbox}[colback=tealclair,colframe=teal,boxrule=1pt,
      left=8pt,right=8pt,top=5pt,bottom=5pt,arc=2pt]
    \textbf{\large Exercice \theexo} \;---\; \textbf{#1}
  \end{tcolorbox}%
  \vspace{4pt}%
  \nopagebreak
}

\newcommand{\partiecours}[1]{%
  \vspace{14pt}%
  {\color{teal}\hrule height 2pt}%
  \vspace{4pt}%
  {\Large\bfseries\color{teal} #1}%
  \vspace{2pt}%
  {\color{teal}\hrule height 0.6pt}%
  \vspace{8pt}%
}

\newtcolorbox{rappel}[1][Rappel]{
  colback=white,colframe=gray!60,boxrule=0.8pt,
  fonttitle=\bfseries,title=#1,coltitle=black,colbacktitle=grisclair,
  left=6pt,right=6pt}

\newtcolorbox{indication}{
  colback=white,colframe=orangefonce!60,boxrule=0.8pt,
  left=6pt,right=6pt,top=4pt,bottom=4pt}

\newcommand{\reponse}{\vspace{4pt}\noindent{\color{teal}\textbf{Solution.}}\;}

\setlength{\headheight}{14pt}
\pagestyle{fancy}
\fancyhf{}
\fancyfoot[C]{\thepage}
\renewcommand{\headrulewidth}{0.4pt}
\renewcommand{\footrulewidth}{0pt}

\fancyhead[L]{\small\color{teal} TD --- Algorithmes de tri}
\fancyhead[R]{\small\color{teal} Énoncé}

\begin{document}

\begin{center}
  {\color{teal}\hrule height 2.5pt}
  \vspace{10pt}
  {\Huge\bfseries\color{teal} Travaux dirigés}\\[6pt]
  {\LARGE\bfseries Les algorithmes de tri}\\[10pt]
  {\large Complexité \;$\bullet$\; Correction \;$\bullet$\; Terminaison}\\[8pt]
  {\normalsize CPGE --- filières MP / PSI \quad$|$\quad Programme officiel marocain 2023}\\[6pt]
  \vspace{6pt}
  {\color{teal}\hrule height 1pt}
\end{center}

\vspace{6pt}

\begin{rappel}[Conventions]
\begin{itemize}[leftmargin=*,itemsep=2pt]
  \item $t$ (ou $L$) désigne un tableau (une liste Python) de $n$ éléments deux à deux
        comparables ; on note $t[g..d]$ la tranche d'indices $g\le k\le d$.
  \item L'\emph{opération élémentaire} de référence est la \textbf{comparaison} entre
        deux éléments ; on compte séparément les \textbf{échanges} et les
        \textbf{décalages} lorsque c'est pertinent.
  \item Une \textbf{inversion} de $t$ est un couple $(i,j)$ tel que $i<j$ et $t[i]>t[j]$.
        On note $I(t)$ leur nombre ; $0\le I(t)\le \binom n2$.
  \item $H_n=\sum_{k=1}^n \frac1k = \ln n + \gamma + o(1)$ ; $\log$ est en base $2$.
  \item Formule de Stirling : $n!\sim\sqrt{2\pi n}\left(\frac ne\right)^{n}$ ;
        en particulier $n!\ge \left(\frac ne\right)^{n}$.
\end{itemize}
\end{rappel}

\begin{rappel}[Méthode --- les trois preuves attendues]
Pour chaque algorithme on demande systématiquement :
\begin{description}[leftmargin=1.4cm,style=nextline,itemsep=2pt]
  \item[Terminaison] un \emph{variant} : quantité entière $\ge0$ strictement
        décroissante à chaque tour (boucle) ; ou une \emph{fonction de mesure} à valeurs
        dans $\mathbb N$ strictement décroissante le long des appels (récursion).
  \item[Correction] un \emph{invariant} de boucle : propriété vraie avant le premier
        tour (initialisation), conservée par un tour (hérédité), et qui donne le résultat
        voulu à la sortie (terminaison de la preuve) ; ou une récurrence bien fondée.
  \item[Complexité] le nombre d'opérations dans le meilleur cas, le pire cas, et
        lorsqu'on le demande en moyenne, avec la classe $\Theta$ finale.
\end{description}
\end{rappel}

% =====================================================================
\partiecours{Partie I --- Les tris quadratiques}
% =====================================================================

\exo{Tri par sélection : comptage exact et invariant}
\begin{lstlisting}
def tri_selection(t):
    n = len(t)
    for i in range(n - 1):
        imin = i
        for j in range(i + 1, n):
            if t[j] < t[imin]:
                imin = j
        t[i], t[imin] = t[imin], t[i]
    return t
\end{lstlisting}

\begin{enumerate}[label=\arabic*.,leftmargin=*]
  \item Donner le nombre \emph{exact} de comparaisons \texttt{t[j] < t[imin]} effectuées,
        en fonction de $n$. Dépend-il du contenu du tableau ?
  \item Donner le nombre exact d'échanges. En déduire la complexité en $\Theta$
        dans le meilleur et dans le pire des cas.
  \item \textbf{Terminaison.} Justifier la terminaison des deux boucles.
  \item \textbf{Correction.} Montrer l'invariant de la boucle externe :
        \begin{quote}
          \emph{avant l'itération $i$, la tranche $t[0..i-1]$ est triée par ordre
          croissant et contient les $i$ plus petits éléments du tableau initial.}
        \end{quote}
        Conclure que l'algorithme trie $t$.
  \item Pourquoi la boucle externe s'arrête-t-elle à $i=n-2$ et non $i=n-1$ ?
\end{enumerate}

% ---------------------------------------------------------------------
\exo{Tri par insertion : le rôle des inversions}
\begin{lstlisting}
def tri_insertion(t):
    for i in range(1, len(t)):
        cle = t[i]
        j = i - 1
        while j >= 0 and t[j] > cle:
            t[j + 1] = t[j]
            j -= 1
        t[j + 1] = cle
    return t
\end{lstlisting}

\begin{enumerate}[label=\arabic*.,leftmargin=*]
  \item Dérouler l'algorithme sur $t=[5,2,4,6,1,3]$ : donner l'état du tableau après
        chaque itération de la boucle \texttt{for}.
  \item \textbf{Meilleur cas.} Quel est le nombre de comparaisons si $t$ est déjà trié ?
        En déduire la complexité.
  \item \textbf{Pire cas.} Même question pour un tableau trié en ordre \emph{décroissant}
        strict. Donner le nombre exact de comparaisons et de décalages.
  \item \textbf{Cas général.} Montrer que le nombre total de décalages
        \texttt{t[j+1] = t[j]} est exactement $I(t)$, le nombre d'inversions du tableau
        initial. \emph{Indication :} montrer que chaque décalage supprime exactement une
        inversion, et qu'aucune opération n'en crée.
  \item En déduire que la complexité du tri par insertion est
        $\Theta\big(n + I(t)\big)$. Pourquoi dit-on que ce tri est
        « adapté aux tableaux presque triés » ? Que vaut la complexité si chaque élément
        est à distance au plus $k$ de sa position finale ?
\end{enumerate}

% ---------------------------------------------------------------------
\exo{Tri par insertion : preuves complètes}
On reprend le code de l'exercice~2.

\begin{enumerate}[label=\arabic*.,leftmargin=*]
  \item \textbf{Terminaison de la boucle \texttt{while}.} Proposer un variant et rédiger
        la preuve. Pourquoi l'accès \texttt{t[j]} ne provoque-t-il jamais d'erreur
        d'indice ? (Le rôle de l'évaluation paresseuse du \texttt{and} en Python.)
  \item \textbf{Correction de la boucle \texttt{while}.} Montrer l'invariant :
        \begin{quote}
          \emph{$t[j+2..i]$ contient, triés, les éléments de $t[j+1..i-1]$ initiaux, tous
          strictement supérieurs à \texttt{cle} ; et $t[0..j]$ est inchangé et trié.}
        \end{quote}
  \item \textbf{Correction de la boucle \texttt{for}.} Énoncer et démontrer l'invariant
        « $t[0..i-1]$ est trié et contient les $i$ premiers éléments initiaux ».
        Conclure.
  \item Justifier que le tri par insertion est un tri \textbf{en place}.
\end{enumerate}

% ---------------------------------------------------------------------
\exo{Tri à bulles et arrêt anticipé}
\begin{lstlisting}
def tri_bulles(t):
    n = len(t)
    for i in range(n - 1):
        echange = False
        for j in range(n - 1 - i):
            if t[j] > t[j + 1]:
                t[j], t[j + 1] = t[j + 1], t[j]
                echange = True
        if not echange:
            return t
    return t
\end{lstlisting}

\begin{enumerate}[label=\arabic*.,leftmargin=*]
  \item \textbf{Correction.} Montrer l'invariant de la boucle interne : \emph{après le
        tour d'indice $j$, $t[j+1]$ est le maximum de $t[0..j+1]$}. En déduire
        l'invariant de la boucle externe : \emph{après l'itération $i$, la tranche
        $t[n-1-i..n-1]$ est triée et contient les $i+1$ plus grands éléments}.
  \item \textbf{Terminaison.} Donner un variant pour chacune des deux boucles.
  \item Donner le nombre de comparaisons dans le pire des cas (exact), et la complexité
        associée.
  \item Grâce au drapeau \texttt{echange}, quelle est la complexité dans le meilleur des
        cas ? Donner une famille de tableaux l'atteignant.
  \item Montrer que le nombre d'échanges effectués est exactement $I(t)$.
        Comparer avec le tri par insertion : lequel effectue le moins de comparaisons
        sur un tableau presque trié ?
\end{enumerate}

% ---------------------------------------------------------------------
\exo{Tri en place, tri stable}
\begin{rappel}[Définitions]
Un tri est \textbf{en place} s'il n'utilise qu'un nombre $\Theta(1)$ d'emplacements
mémoire auxiliaires, indépendant de $n$ (la pile des appels récursifs comptant dans
l'espace utilisé).
Un tri est \textbf{stable} si, pour tout couple d'indices $i<j$ tel que
$t[i]$ et $t[j]$ sont de clés \emph{égales}, l'élément initialement en position $i$ se
retrouve, après le tri, à une position strictement inférieure à celle de l'élément
initialement en position $j$.
\end{rappel}

\begin{enumerate}[label=\arabic*.,leftmargin=*]
  \item Pourquoi la stabilité n'a-t-elle de sens que lorsque l'on trie des
        \emph{enregistrements} selon une \emph{clé} ? Donner une situation concrète où
        elle est indispensable (tri d'une liste d'élèves par note, puis par classe).
  \item Montrer que le tri par sélection donné en exercice~1 n'est \textbf{pas stable} :
        exhiber un tableau de $3$ éléments (clés notées $2_a, 2_b, 1$) qui le prouve.
  \item Montrer que le tri par insertion \textbf{est stable}. Quel est le rôle exact
        du test \texttt{t[j] > cle} (et non \texttt{>=}) ?
  \item Le tri à bulles est-il stable ? Justifier.
  \item Compléter le tableau suivant (en place ? stable ?) pour les quatre tris
        rencontrés jusqu'ici, et le compléter à la fin du TD pour le tri fusion et le
        tri rapide.
\end{enumerate}

% =====================================================================
\partiecours{Partie II --- Le tri par fusion}
% =====================================================================

\exo{La fonction \texttt{fusion} : spécification, invariant, variant}
\begin{lstlisting}
def fusion(A, B):
    C = []
    i, j = 0, 0
    while i < len(A) and j < len(B):
        if A[i] <= B[j]:
            C.append(A[i]); i += 1
        else:
            C.append(B[j]); j += 1
    return C + A[i:] + B[j:]
\end{lstlisting}

\begin{enumerate}[label=\arabic*.,leftmargin=*]
  \item \textbf{Spécification.} Énoncer précisément : préconditions sur \texttt{A} et
        \texttt{B}, postcondition sur la valeur renvoyée.
  \item \textbf{Terminaison.} Donner un variant de la boucle \texttt{while} et rédiger
        la preuve.
  \item \textbf{Correction.} Démontrer l'invariant :
        \begin{quote}
          \emph{$C$ est trié, $C$ est formé exactement des éléments de $A[0..i-1]$ et
          $B[0..j-1]$, et tout élément de $C$ est inférieur ou égal à tout élément
          restant de $A[i..]$ et de $B[j..]$.}
        \end{quote}
        En déduire la correction de \texttt{fusion}.
  \item \textbf{Complexité.} On note $p=|A|$ et $q=|B|$. Montrer que le nombre de
        comparaisons $c$ vérifie
        \[ \min(p,q) \;\le\; c \;\le\; p+q-1 .\]
        Donner un exemple atteignant chacune des deux bornes. Conclure que
        \texttt{fusion} est en $\Theta(p+q)$.
  \item Quelle est la complexité \emph{spatiale} de \texttt{fusion} ?
\end{enumerate}

% ---------------------------------------------------------------------
\exo{\texttt{tri\_fusion} : terminaison et correction}
\begin{lstlisting}
def tri_fusion(L):
    if len(L) <= 1:
        return L
    m = len(L) // 2
    gauche = tri_fusion(L[:m])
    droite = tri_fusion(L[m:])
    return fusion(gauche, droite)
\end{lstlisting}

\begin{enumerate}[label=\arabic*.,leftmargin=*]
  \item \textbf{Terminaison.} Proposer une fonction de mesure et montrer qu'elle
        décroît strictement à chaque appel récursif. On vérifiera soigneusement que
        $0<m<|L|$ dès que $|L|\ge2$ --- que se passerait-il si l'on avait écrit
        \texttt{m = 0} ou \texttt{m = len(L)} ?
  \item \textbf{Correction.} Démontrer par récurrence forte sur $n=|L|$ que
        \texttt{tri\_fusion(L)} renvoie une permutation triée de \texttt{L}.
  \item \texttt{tri\_fusion} est-il \textbf{en place} ? Évaluer l'espace mémoire total
        utilisé (tranches créées par \texttt{L[:m]} et \texttt{L[m:]}, pile d'appels).
  \item Dérouler l'arbre des appels sur $L=[5,2,4,6,1,3]$ : dessiner l'arbre de
        découpage puis les fusions successives.
\end{enumerate}

% ---------------------------------------------------------------------
\exo{Complexité du tri fusion}
On note $C(n)$ le nombre de comparaisons dans le \emph{pire} des cas pour $n$ éléments.

\begin{enumerate}[label=\arabic*.,leftmargin=*]
  \item Établir la récurrence
        $\;C(n)=C(\lfloor n/2\rfloor)+C(\lceil n/2\rceil)+ (n-1)$ pour $n\ge2$,
        avec $C(1)=0$.
  \item On suppose $n=2^{k}$. Résoudre exactement la récurrence
        $C(2^{k}) = 2\,C(2^{k-1}) + 2^{k}-1$ et montrer que
        \[
          C(n) = n\log_2 n - n + 1 .
        \]
        Vérifier pour $n=2$ et $n=4$.
  \item En déduire que le tri fusion est en $\Theta(n\log n)$ dans le pire des cas.
        Que vaut la complexité dans le meilleur des cas ? Le tri fusion est-il sensible
        à l'ordre initial des données ?
  \item \textbf{Cas général.} Montrer par récurrence forte que
        $C(n)\le n\lceil\log_2 n\rceil$ pour tout $n\ge1$.
  \item Comparer numériquement $n\log_2 n$ et $n^2/2$ pour $n=10^{3}$ et $n=10^{6}$.
\end{enumerate}

% ---------------------------------------------------------------------
\exo{Stabilité du tri fusion}
\begin{enumerate}[label=\arabic*.,leftmargin=*]
  \item Montrer que \texttt{tri\_fusion} est stable. On identifiera précisément
        l'endroit du code qui assure cette propriété.
  \item On remplace le test \texttt{A[i] <= B[j]} par \texttt{A[i] < B[j]}.
        L'algorithme reste-t-il correct ? Reste-t-il stable ? Donner un contre-exemple
        explicite avec des clés égales.
  \item On remplace \texttt{m = len(L) // 2} par \texttt{m = 1}. L'algorithme
        reste-t-il correct ? Quelle est alors sa complexité ? Quel tri connu
        retrouve-t-on ?
\end{enumerate}

% =====================================================================
\partiecours{Partie III --- Le tri rapide}
% =====================================================================

\exo{\texttt{tri\_rapide} : preuves et cas extrêmes}
\begin{lstlisting}
def tri_rapide(L):
    if len(L) <= 1:
        return L
    pivot = L[0]
    petits  = [x for x in L[1:] if x <= pivot]
    grands  = [x for x in L[1:] if x >  pivot]
    return tri_rapide(petits) + [pivot] + tri_rapide(grands)
\end{lstlisting}

\begin{enumerate}[label=\arabic*.,leftmargin=*]
  \item \textbf{Terminaison.} Proposer une fonction de mesure. Justifier soigneusement
        que $|\texttt{petits}|<|L|$ \emph{et} $|\texttt{grands}|<|L|$ --- où intervient
        le fait que le pivot est retiré de la liste ?
  \item \textbf{Correction.} Démontrer par récurrence forte sur $|L|$ que
        \texttt{tri\_rapide(L)} renvoie une permutation triée de \texttt{L}.
        On précisera pourquoi la concaténation
        \texttt{tri\_rapide(petits) + [pivot] + tri\_rapide(grands)} est bien triée.
  \item \textbf{Meilleur cas.} On suppose qu'à chaque appel le pivot est un élément
        médian, c'est-à-dire $|\texttt{petits}|=\lfloor (n-1)/2\rfloor$. Écrire la
        récurrence du nombre de comparaisons et la résoudre. Conclusion.
  \item \textbf{Pire cas.} Montrer que si $L$ est déjà trié par ordre croissant, alors
        \texttt{petits} est vide à chaque appel. Écrire la récurrence
        correspondante, la résoudre exactement et conclure.
        \emph{Le tri rapide est donc quadratique sur une entrée déjà triée} :
        commenter ce paradoxe apparent.
  \item Déterminer la profondeur maximale de la pile d'appels dans le meilleur et dans
        le pire des cas. Quel risque pratique cela représente-t-il pour $n=10^{6}$ ?
\end{enumerate}

% ---------------------------------------------------------------------
\exo{Complexité en moyenne du tri rapide}
On suppose que $L$ est une permutation aléatoire uniforme de $n$ éléments deux à deux
distincts. On note $C(n)$ l'espérance du nombre de comparaisons, avec $C(0)=C(1)=0$.

\begin{enumerate}[label=\arabic*.,leftmargin=*]
  \item Justifier que le rang du pivot dans $L$ suit la loi uniforme sur
        $\{0,1,\dots,n-1\}$, puis établir
        \[
          C(n) = (n-1) + \frac1n\sum_{k=0}^{n-1}\big(C(k)+C(n-1-k)\big)
               = (n-1) + \frac2n\sum_{k=0}^{n-1}C(k).
        \]
  \item Multiplier par $n$, écrire la même relation au rang $n-1$ et soustraire pour
        obtenir
        \[
          n\,C(n) = (n+1)\,C(n-1) + 2(n-1).
        \]
  \item Poser $D(n)=\dfrac{C(n)}{n+1}$ et montrer que
        $D(n)=D(n-1)+\dfrac{2(n-1)}{n(n+1)}$.
        Décomposer $\dfrac{2(n-1)}{n(n+1)}$ en éléments simples.
  \item En déduire $D(n) = 2H_n + \dfrac{4}{n+1} - 4$, puis
        \[
          \boxed{\;C(n) = 2(n+1)H_n - 4n\;}
        \]
        Vérifier la formule pour $n=2$ et $n=3$.
  \item En déduire $C(n)\sim 2n\ln n = 2\ln 2\cdot n\log_2 n \approx 1{,}39\,n\log_2 n$.
        Comparer au tri fusion ($\approx 1\cdot n\log_2 n$) : de combien de pour cent le
        tri rapide effectue-t-il plus de comparaisons en moyenne ?
  \item Le tri rapide est pourtant réputé plus rapide en pratique. Proposer deux raisons.
\end{enumerate}

% ---------------------------------------------------------------------
\exo{Partition en place (schéma de Lomuto)}
La version par compréhension de l'exercice~10 n'est \emph{pas} en place. On lui
substitue la partition suivante, qui trie la tranche $t[g..d]$ en utilisant $t[d]$
comme pivot.

\begin{lstlisting}
def partition(t, g, d):
    p = t[d]
    i = g - 1
    for j in range(g, d):
        if t[j] <= p:
            i += 1
            t[i], t[j] = t[j], t[i]
    t[i+1], t[d] = t[d], t[i+1]
    return i + 1

def tri_rapide_place(t, g, d):
    if g < d:
        m = partition(t, g, d)
        tri_rapide_place(t, g, m - 1)
        tri_rapide_place(t, m + 1, d)
\end{lstlisting}

\begin{enumerate}[label=\arabic*.,leftmargin=*]
  \item Dérouler \texttt{partition} sur $t=[2,8,7,1,3,5,6,4]$ avec $g=0$, $d=7$ :
        donner l'état du tableau et les valeurs de $i,j$ à chaque tour.
  \item \textbf{Correction.} Démontrer l'invariant de la boucle \texttt{for} :
        \begin{quote}
          \emph{pour $g\le k\le i$ : $t[k]\le p$ ; \quad pour $i+1\le k\le j-1$ :
          $t[k]>p$ ; \quad et $t[d]=p$.}
        \end{quote}
        En déduire qu'après l'échange final, l'indice $m=i+1$ renvoyé vérifie :
        $t[m]=p$, tous les éléments de $t[g..m-1]$ sont $\le p$ et tous ceux de
        $t[m+1..d]$ sont $>p$.
  \item \textbf{Terminaison.} Justifier la terminaison de \texttt{partition}, puis celle
        de \texttt{tri\_rapide\_place} (fonction de mesure sur $d-g$).
  \item Donner le nombre exact de comparaisons effectuées par un appel à
        \texttt{partition(t,g,d)}.
  \item Quelle est la complexité \emph{spatiale} de \texttt{tri\_rapide\_place} dans le
        meilleur et le pire des cas ? Cet algorithme est-il en place au sens strict ?
  \item Cette version est-elle stable ? Justifier par un contre-exemple.
\end{enumerate}

% ---------------------------------------------------------------------
\exo{Choix du pivot et entrées dégénérées}
\begin{enumerate}[label=\arabic*.,leftmargin=*]
  \item Soit $t$ un tableau dont tous les éléments sont \emph{égaux}. Montrer qu'avec la
        partition de Lomuto (test \texttt{t[j] <= p}), la fonction \texttt{partition}
        renvoie toujours $m=d$. En déduire la complexité de
        \texttt{tri\_rapide\_place} sur un tel tableau. Même question pour la version
        par compréhension de l'exercice~10.
  \item On choisit désormais le pivot \emph{au hasard} parmi $t[g..d]$ (en l'échangeant
        d'abord avec $t[d]$). Le pire des cas $\Theta(n^2)$ existe-t-il encore ?
        Qu'est-ce qui change exactement ?
  \item \textbf{Pivot médian de trois.} On prend pour pivot la médiane de
        $t[g]$, $t[(g+d)//2]$, $t[d]$. Montrer que sur un tableau déjà trié, la
        partition devient équilibrée. Existe-t-il encore des entrées quadratiques ?
  \item \textbf{Partition ternaire (ouverture).} On partitionne en trois zones
        $(<p)$, $(=p)$, $(>p)$ et l'on ne rappelle récursivement que sur les zones
        extrêmes. Quelle est alors la complexité sur un tableau constant ?
        Sur un tableau ne prenant que $d$ valeurs distinctes ?
\end{enumerate}

% =====================================================================
\partiecours{Partie IV --- Borne inférieure, applications, synthèse}
% =====================================================================

\exo{La borne $\Omega(n\log n)$ : arbre de décision}
On appelle \emph{tri par comparaisons} tout algorithme dont les seules opérations sur
les éléments sont des comparaisons deux à deux. Un tel algorithme, exécuté sur $n$
éléments deux à deux distincts, se représente par un \textbf{arbre de décision} :
chaque nœud interne est une comparaison $t[i]<t[j]$ (deux fils : oui / non), chaque
feuille est une permutation renvoyée.

\begin{enumerate}[label=\arabic*.,leftmargin=*]
  \item Dessiner l'arbre de décision du tri par insertion pour $n=3$.
        Combien a-t-il de feuilles ? Quelle est sa hauteur ?
  \item Justifier qu'un arbre de décision correct pour $n$ éléments possède
        \emph{au moins} $n!$ feuilles.
  \item Montrer qu'un arbre binaire de hauteur $h$ a au plus $2^{h}$ feuilles.
        En déduire $h \ge \log_2(n!)$.
  \item À l'aide de $n!\ge (n/e)^{n}$, montrer que
        \[
          h \;\ge\; n\log_2 n - n\log_2 e \;=\; n\log_2 n - 1{,}4427\,n .
        \]
        Conclure : tout tri par comparaisons effectue $\Omega(n\log n)$ comparaisons
        dans le pire des cas.
  \item Le tri fusion est-il optimal au sens de cette borne ? Et le tri rapide ?
  \item Pourquoi cette borne ne s'applique-t-elle pas à un tri qui utiliserait les
        éléments comme des \emph{indices} de tableau (tri par comptage) ?
\end{enumerate}

% ---------------------------------------------------------------------
\exo{Trier pour résoudre}
Dans tous les cas on comparera l'approche naïve et l'approche « on trie d'abord ».

\begin{enumerate}[label=\arabic*.,leftmargin=*]
  \item \textbf{Doublons.} Écrire une fonction qui décide si un tableau de $n$ entiers
        contient deux éléments égaux. Donner la complexité de la version naïve à deux
        boucles, puis celle de la version par tri. Conclure.
  \item \textbf{Paire de somme donnée.} Étant donnés $t$ et $s$, décider s'il existe
        $i\ne j$ avec $t[i]+t[j]=s$. Proposer un algorithme en $\Theta(n\log n)$ et
        justifier sa complexité (on pourra réutiliser la méthode des deux indices).
  \item \textbf{Les $k$ plus grands éléments.} Comparer :
        (i) $k$ tours de tri par sélection, de complexité $\Theta(kn)$ ;
        (ii) un tri complet suivi d'une lecture, en $\Theta(n\log n)$.
        Pour quelles valeurs de $k$ la première approche est-elle préférable ?
  \item \textbf{Élément majoritaire.} Un élément est majoritaire s'il apparaît
        strictement plus de $n/2$ fois. Proposer un algorithme en $\Theta(n\log n)$
        fondé sur le tri, et en prouver la correction.
  \item \textbf{Médiane.} Quelle est la complexité du calcul de la médiane par tri ?
\end{enumerate}

% ---------------------------------------------------------------------
\exo{Synthèse et questions de cours}
\begin{enumerate}[label=\arabic*.,leftmargin=*]
  \item Compléter le tableau récapitulatif :
        \begin{center}\small
        \begin{tabular}{|l|c|c|c|c|c|}
        \hline
        \textbf{Tri} & \textbf{Meilleur} & \textbf{Moyenne} & \textbf{Pire} &
        \textbf{En place} & \textbf{Stable} \\\hline
        Sélection  & & & & & \\\hline
        Insertion  & & & & & \\\hline
        Bulles     & & & & & \\\hline
        Fusion     & & & & & \\\hline
        Rapide (Lomuto) & & & & & \\\hline
        \end{tabular}
        \end{center}
  \item \textbf{Vrai ou faux}, en justifiant systématiquement :
        \begin{enumerate}[label=(\alph*),itemsep=1pt]
          \item Un tri en place est nécessairement stable.
          \item Un algorithme en $\Theta(n\log n)$ est toujours plus rapide qu'un
                algorithme en $\Theta(n^2)$.
          \item Le tri rapide est en $\Theta(n\log n)$.
          \item Si $I(t)=O(n)$, le tri par insertion est linéaire.
          \item Aucun tri ne peut faire moins de $n\log_2 n$ comparaisons.
        \end{enumerate}
  \item \textbf{Tri hybride.} On modifie \texttt{tri\_fusion} : lorsque
        $|L|\le S$ (seuil fixé), on trie $L$ par insertion au lieu de découper.
        \begin{enumerate}[label=(\alph*),itemsep=1pt]
          \item Justifier que l'algorithme reste correct.
          \item Montrer que la complexité dans le pire des cas est
                $\Theta\!\big(nS + n\log(n/S)\big)$.
          \item Pour quel ordre de grandeur de $S$ cette expression reste-t-elle
                $\Theta(n\log n)$ ? Que se passe-t-il si $S=\Theta(n)$ ?
          \item Pourquoi ce choix améliore-t-il les performances réelles alors que
                la classe de complexité est inchangée ?
        \end{enumerate}
\end{enumerate}

\vspace{14pt}
{\color{teal}\hrule height 1pt}
\vspace{6pt}
\begin{center}
\small\itshape Fin de l'énoncé --- 16 exercices.
\end{center}

\end{document}
